problem:  
Let \((X,d,\mathfrak{m})\) be a compact, essentially non‑branching metric measure space satisfying the curvature–dimension condition \(\mathrm{CD}(K,N)\) for some \(K\in\mathbb{R}\), \(1<N<\infty\).  
For bounded Borel functions \(f\in L^\infty(X,\mathfrak{m})\) and \(t_0>0\), set \(\rho_0 := P_{t_0}f\), where \(P_t\) is the heat semigroup on \((X,d,\mathfrak{m})\); thus \(\mu_0=\rho_0\mathfrak{m}\) is an absolutely continuous probability measure with \(L^\infty\) density.  

For each Lipschitz function \(\phi\) and each regularization parameter \(\varepsilon>0\), denote by \(\mu_1^{\varepsilon,\phi}\) the terminal marginal of the entropic interpolation starting from \(\mu_0\) and generated by the Schrödinger bridge with initial potential \(\phi\) (i.e. the marginal at time \(t=1\) of the minimizer of the entropy-regularized cost with reference measure \(e^{-d^2/(2\varepsilon)}\mathfrak{m}^{\otimes2}\)).  

Define the **entropic second‑order entropy increment** functional
\[
\mathscr{E}_\varepsilon[\mu_0,\phi]
:=
\frac{1}{\varepsilon^2}\left(
\mathrm{Ent}_{\mathfrak{m}}(\mu_1^{\varepsilon,\phi})
- \mathrm{Ent}_{\mathfrak{m}}(\mu_0)
- \frac{\varepsilon}{2}\!\int_X\! |\nabla^- \phi|^2\,d\mu_0
\right),
\]
where \(|\nabla^- \phi|\) denotes the metric slope of \(\phi\).  

Assume that for all admissible \(\mu_0\) and every pair of Lipschitz potentials \(\phi_1,\phi_2\), the limit  
\[
\mathcal{R}(\phi_1,\phi_2)
:= \lim_{\varepsilon\to0}\frac{1}{2}\Big(\mathscr{E}_\varepsilon[\mu_0,\phi_1+\phi_2]
- \mathscr{E}_\varepsilon[\mu_0,\phi_1]
- \mathscr{E}_\varepsilon[\mu_0,\phi_2]\Big)
\]
exists finitely, is symmetric, and independent of the initial measure \(\mu_0\).  

**Question:**  
Does the existence and base‑measure independence of this limiting bilinear form \(\mathcal{R}\) imply that \((X,d,\mathfrak{m})\) admits a tensorial “entropic Ricci curvature” compatible with the metric, i.e. that it must in fact be an \(\mathrm{RCD}(K,N)\) space?  
Equivalently, can the infinitesimal second‑order expansion of the entropy along entropic interpolations serve to *define* a genuine Ricci curvature tensor synthetically, recovering Riemannian structure from purely entropic data?  

Why is it a "good" problem:  
This version is both conceptually and mathematically well‑posed within existing synthetic geometry.  
All objects—heat flow regularization, metric slopes, entropic interpolations, and entropy—are defined for CD spaces, avoiding ill‑defined deterministic transport maps while retaining the profound geometric content.  

1. **Conceptual novelty and depth:**  
   The question proposes a fundamentally new analytic object, the *entropic curvature form* \(\mathcal{R}\), built from second‑order asymptotics of entropy under entropic optimal transport.  
   It asks whether the very existence and density‑independence of this bilinear form force the space to be infinitesimally Riemannian.  

2. **Cross‑disciplinary synthesis:**  
   The problem bridges stochastic analysis (Schrödinger bridge theory), calculus in metric measure spaces, and curvature‑dimension geometry. It draws a deep analogy with the Bakry–Émery \(\Gamma_2\) operator, but based solely on entropic transports rather than diffusion semigroups.  

3. **Logical soundness and simplicity:**  
   Each component—metric slope, entropy, entropic interpolations—is rigorously defined for any CD space, so the problem is precise, self‑consistent, and yet surprisingly compact.  
   The statement is simple (“Does second‑order entropy expansion determine curvature?”), but its resolution would require the full machinery of Wasserstein geometry and stochastic regularization analysis.  

4. **Potential for foundational impact:**  
   If true, this would yield an operator‑based analytic way to reconstruct Ricci curvature tensors in nonsmooth settings, unifying entropy, curvature, and information geometry.  
   Conversely, counterexamples would pinpoint the exact analytical gaps preventing synthetic Ricci tensors from being defined solely through entropy transport, delineating new frontiers in metric measure geometry.  

Thus, the problem retains the elegant concision of its motivation and realizes a mathematically coherent, profoundly exploratory question that sits at the heart of synthetic Ricci curvature via optimal transport.